\documentclass[aspectratio=169]{beamer}

% Use Sorbonne theme
\usepackage{sorbonne_theme}

\DeclareMathOperator{\plim}{plim}

% Title page information
\title{Econometrics Part 5}
\subtitle{Endogeneity}
\author{Jose Angel Garcia Sanchez}
\institute{Université Paris 1 Panthéon-Sorbonne}
\date{} % left empty on purpose: beamer would otherwise print today's date

% Set the logo
\sorbonnelogo{imagen.png}

\begin{document}

% Title page (without logo in corner)
{
\setbeamertemplate{logo}{}
\begin{frame}
\titlepage
\end{frame}
}

% Table of contents
\begin{frame}{Outline}
\tableofcontents
\end{frame}

% Section 1: Framework
\section{Framework}

\begin{frame}{Considering random explanatory variables}
\begin{itemize}
    \item In the basic framework, explanatory variables $X$ are usually considered not random
    \item This amounts to reasoning conditionally on $X$, i.e.\ as if the $X$ were fixed
    \item An important hypothesis was that $X$ and $u$ were independent, which is always the case if $X$ is fixed
    \item But the $X$'s can be considered as random variables
    \item All the previous results still hold (as long as $E(u|X) = 0$), because they were obtained conditionally on $X$
\end{itemize}
\end{frame}

\begin{frame}{Assumptions for estimation}
\begin{enumerate}
    \item $E(u) = 0$
    \item $X$ is random and $E(u|X) = 0$: in particular, $\forall t, t'$, $X_t$ is uncorrelated with $u_{t'}$
    \item $\operatorname{rank}(X) = k$
    \item $E(uu') = \sigma^2I_N$
    \item Error terms are iid $(0, \sigma^2)$
    \item $\plim(X'X/N) = V_X$, which is a positive definite matrix (when $N$ goes to infinity, the $X$ variables always keep some variance)
\end{enumerate}
\end{frame}

\begin{frame}{Properties of OLS: unbiasedness}
\begin{itemize}
    \item $\hat{b}_{ols}$ is still unbiased
    \item We know that $\hat{b}_{ols} = (X'X)^{-1}X'y = b + (X'X)^{-1}X'u$
    \item So $E(\hat{b}_{ols}) = b + E[(X'X)^{-1}X'u]$
    \item We get: $E[(X'X)^{-1}X'u] = E_X\big[E_u[(X'X)^{-1}X'u \mid X]\big]$\\
    $= E_X\big[(X'X)^{-1}X'E_u[u \mid X]\big]$ and $E_u[u \mid X] = 0$
    \item In the same way, we get an expression for the variance of $\hat{b}_{ols}$:
    \item $V(\hat{b}_{ols}) = \sigma^2E[(X'X)^{-1}]$
\end{itemize}
\end{frame}

\begin{frame}{Properties of OLS: consistency}
$\hat{b}_{ols}$ is still consistent:

\[\plim(\hat{b}_{ols}) = b + \left[\plim\left(\frac{X'X}{N}\right)\right]^{-1}\plim\left(\frac{X'u}{N}\right) = b + V_X^{-1} \cdot 0 = b\]

We will see on the next slide that the asymptotic distribution of $\hat{b}$ is a normal, as expected.
\end{frame}

\begin{frame}{Asymptotics: OLS estimators}
We wish now to enlarge our framework to the case where $u$'s are only iid (and not necessarily normal) and where $X$'s can be random (and not necessarily fixed). In what follows, we will thus use the following asymptotic results, obtained under the usual model hypotheses:

\[\plim(\hat{b}_{ols}) = b + \left[\plim\left(\frac{X'X}{N}\right)\right]^{-1}\plim\left(\frac{X'u}{N}\right) = b + V_X^{-1} \cdot 0 = b\]

\[\plim(\hat{\sigma}_{ols}^2) = \sigma^2\]

\[\sqrt{N}(\hat{b}_{ols} - b) \xrightarrow{d} \mathcal{N}(0, \sigma^2V_X^{-1})\]
\end{frame}

\begin{frame}{Asymptotic tests}
Tests of $r$ linear constraints ($C$ is of size $(r,k)$) are run using the following formula:

\[f = \frac{(C\hat{b}_{ols} - Cb)'(C(X'X)^{-1}C')^{-1}(C\hat{b}_{ols} - Cb)}{\sigma^2} \xrightarrow{d} \chi_r^2\]

Where, as before in the finite sample case, $\sigma^2$ is unknown. But since by assumption we work with a very large sample, $\hat{\sigma}_{ols}^2$ is a consistent estimate of $\sigma^2$. We can thus replace $\sigma^2$ by its consistent estimate, without having to replace the $\chi^2$ with a Fisher.
\end{frame}

\begin{frame}{General properties}
If we assume that the $X_i$ are iid and the $u$'s are normal, then:
\begin{itemize}
    \item OLS estimators are consistent, and normal conditionally on $X$
    \item The T- and F-statistics keep their Student and Fisher distributions (their conditional distribution does not depend on $X$): we can run the usual T- and F-tests
    \item If the $u$ are not normal but only iid, we need to make sure the sample is large, and use the asymptotic version of the usual T- and F-tests
    \item In that case, Normal replaces Student and $\chi^2$ replaces Fisher
\end{itemize}
\end{frame}

% Section 2: Non-exogeneity
\section{Non-exogeneity}

\begin{frame}{Non-exogenous explanatory variables}
\begin{itemize}
    \item Let's use the following model: $y_t = X_t'b + u_t$
    \item Assume that $E(X_tu_t) \neq 0$: some explanatory variables are correlated with the current error term
    \item In this case, $\plim(X'u/N) \neq 0$
    \item OLS is biased and inconsistent
\end{itemize}

\[\plim(\hat{b}_{ols}) = b + \left[\plim\left(\frac{X'X}{N}\right)\right]^{-1}\plim\left(\frac{X'u}{N}\right) \neq b\]

Even if we increase the size of the sample, we won't get the right value for $b$.
\end{frame}

\begin{frame}{Exogeneity: definitions}
\begin{itemize}
    \item \textbf{Case 1}: Strict exogeneity: $\forall t, t'$, $E(X_tu_{t'}) = 0$
    \item $X$ and $u$ are never correlated
    \item \textbf{Case 2}: (Contemporaneous) exogeneity: $\forall t$, $E(X_tu_t) = 0$
    \item $X$ and $u$ are uncorrelated for any given time period
    \item \textbf{Case 3}: $X$ predetermined: $\forall t' \geq t$, $E(X_tu_{t'}) = 0$
    \item $X$ uncorrelated with current and future $u$
\end{itemize}

This class focuses on cross-sections, so on the second case.
\end{frame}

\begin{frame}{Omitted variables}
Let's assume the true model is the following:

\[y_t = X_t'b + w_td + v_t\]

If we omit $w_t$ and estimate the following model instead:

\[y_t = X_t'b + u_t\]

\begin{itemize}
    \item $w_t$ is included in $u_t$
    \item If $w_t$ is correlated with $X$, then $E(X_tw_t) \neq 0$ and thus $E(X_tu_t) \neq 0$: OLS is biased
    \item Remark: in the fixed $X$ case, we would get the same result using the Frisch-Waugh theorem
    \item When a variable that is correlated with the $X$'s is omitted, estimators suffer from an \emph{omitted variable bias}
    \item Estimators are inconsistent
\end{itemize}
\end{frame}

\begin{frame}{Variables measured with error}
Let's use a very simple one-variable theoretical model (no constant if variables are centered):

\[y_t^* = x_t^* b + v_t\]

Assume that $x$ is measured with error. To run the estimation, we have access to values $y_t = y_t^*$ and $x_t = x_t^* + e_t$, $e_t$ being a white noise. The estimated model is thus $y_t = x_t b + u_t$, with $u_t = v_t - be_t$. We thus have:

\[E(x_tu_t) = E[(x_t^* + e_t)(v_t - be_t)] = -b\sigma_e^2 \neq 0\]

OLS is thus inconsistent, and biased towards 0 (\textit{attenuation bias}): $\plim(\hat{b}_{ols}) = b\,\sigma_{x^*}^2/(\sigma_{x^*}^2 + \sigma_e^2)$, so the influence of $x$ on $y$ is underestimated.
\end{frame}

\begin{frame}{Simultaneous equations}
Say we have the following system of equations:
\begin{itemize}
    \item $Y_t = a + bX_t + u_t$ (1)
    \item $X_t = Y_t + Z_t$ (2)
    \item Because of equation (1), $Y_t$ is endogenous, and because of equation (2), $X_t$ is endogenous too
    \item At the system level, only $Z_t$ is exogenous
    \item To get back to what we are used to, we should rewrite endogenous variables as functions of exogenous variables only
\end{itemize}

The system can be rewritten (if $b \neq 1$): $Y_t = \frac{a}{1-b} + \frac{b}{1-b}Z_t + \frac{1}{1-b}u_t$ and $X_t = \frac{a}{1-b} + \frac{1}{1-b}Z_t + \frac{1}{1-b}u_t$. Notice that $X_t$ is a function of $u_t$: in equation (1), it is thus correlated with the error term. A basic hypothesis of OLS is violated, and if we estimate (1) without taking into account the information provided by (2), estimates will be inconsistent.
\end{frame}

\begin{frame}{Autocorrelation with lagged dependent variable}
Assume we use the following model:

\[y_t = b_0 + b_1x_t + b_2y_{t-1} + u_t\]

With: $u_t = \rho u_{t-1} + \varepsilon_t$

We can then write:
\begin{itemize}
    \item $y_t = b_0 + b_1x_t + b_2y_{t-1} + \rho u_{t-1} + \varepsilon_t$ (1)
    \item And at the same time: $y_{t-1} = b_0 + b_1x_{t-1} + b_2y_{t-2} + u_{t-1}$ (2)
    \item Equation (2) shows that $y_{t-1}$ depends on $u_{t-1}$, so $y_{t-1}$ is correlated with the error term $\rho u_{t-1} + \varepsilon_t$ in equation (1)
    \item The lagged outcome variable is thus not exogenous in the estimated model
    \item If error terms are not autocorrelated ($\rho = 0$), $y_{t-1}$ is uncorrelated with $u_t$: it is predetermined (case 3)
    \item How to test for this: Durbin's $h$ statistic
\end{itemize}
\end{frame}

\begin{frame}{General issue in policy evaluation}
\begin{itemize}
    \item Say we want to evaluate the impact of a policy on people's wages (e.g.\ a training program)
    \item A model describing the wage outcome is $Y_i = a + bX_i + cP_i + u_i$, where $X$ comprises individual characteristics and $P$ is a dummy variable indicating whether the individual was assigned to the program
    \item If the policy is not \emph{randomized}, i.e.\ if being assigned to the program depends on unobserved individual characteristics, then $P$ is correlated with $u$ and the impact of the policy cannot be estimated consistently
    \item Intuitive reason: those who were assigned to the program are not comparable to those who were not, so the latter cannot be considered as a valid control group for the former
    \item This is called a \emph{selection effect}
\end{itemize}
\end{frame}

\begin{frame}{Consequences}
In the case of non-exogeneity (also called endogeneity) of some explanatory variables:
\begin{itemize}
    \item OLS estimators are biased: on average, we do not get the true value of parameters
    \item OLS estimators are inconsistent: even if we increase the size of the sample, the bias does not go to zero and we will never get the true value of parameters
\end{itemize}

We thus have to use another estimation technique, using auxiliary variables: the instrumental variables technique.
\end{frame}

% Section 3: The IV estimator
\section{The IV estimator}

\begin{frame}{Instrumental variables}
Let's use the following model, where (to make things simple) \emph{all} $X$'s are endogenous, except the constant (which can't be endogenous by definition):

\[y = Xb + u\]

Let's consider a set of variables $Z_{(N,p)}$ different from the $X$'s but including the constant, with the following properties:
\begin{itemize}
    \item $E(Z_tu_t) = 0$: variables $Z$ are exogenous
    \item $\operatorname{rank}(Z) = p$
    \item $\plim(Z'X/N) = V_{ZX}$, with $V_{ZX}$ a non-null matrix of dimension $(p,k)$ and rank $k$
    \item $\plim(Z'Z/N) = V_Z$, with $V_Z$ a finite positive definite matrix of dimension $(p,p)$
\end{itemize}

Meaning: variables $Z$ need to be both exogenous and correlated with $X$, and the number of instruments $p$ must be at least the number of explanatory variables $k$ ($p \geq k$).
\end{frame}

\begin{frame}{The estimator}
The original model is:

\[y = Xb + u\]

\begin{itemize}
    \item Assume we regress $X$ on variables $Z$
    \item We compute their predictions $\tilde{X}$
    \item We use these new variables in the model instead:
\end{itemize}

\[y = \tilde{X}b + u\]

\begin{itemize}
    \item It should work because $\tilde{X}$ is only made up of $Z$ that are exogenous: $\tilde{X}$ is thus exogenous too
\end{itemize}
\end{frame}

\begin{frame}{Intuition (1)}
\begin{itemize}
    \item OLS on the original model $y = Xb + u$ are inconsistent because variables $X$ are correlated with $u$
    \item We can't replace the $X$'s by other variables because the model should remain unchanged
    \item We would like to get rid of this correlation, while keeping the information provided by $X$
    \item Imagine a regression where each variable in $X$ would be explained by the set of variables $Z$, say: $\forall j$, $x_j = Zd + e$
    \item A prediction for $x_j$ would be: $\hat{x}_j = Z\hat{d}$
    \item If this regression is good enough, $\hat{x}_j$ is close to the true $x_j$
    \item And $\hat{x}_j$ is a linear combination of variables $Z$ that are exogenous: it is thus exogenous as well
    \item We can thus replace $x_j$ by $\hat{x}_j$ in the original model
\end{itemize}
\end{frame}

\begin{frame}{Transforming the model}
The original model is:

\[y = Xb + u\]

The new model is:

\[y = \tilde{X}b + u\]

\begin{itemize}
    \item Problem: the new model is not directly equivalent to the original one: only $X$ was modified
    \item How can we make the new model equivalent to the old one?
    \item Notice that $\tilde{X}$ is the prediction of $X$ using the $Z$
    \item So that $\tilde{X} = P_ZX$
    \item We need to multiply every term of the original model by $P_Z$
\end{itemize}
\end{frame}

\begin{frame}{The IV estimator}
\begin{itemize}
    \item Let's call $P_Z$ the orthogonal projection matrix that projects on $L(Z)$: $P_Z = Z(Z'Z)^{-1}Z'$
    \item $P_Z$ is symmetric ($P_Z' = P_Z$) and idempotent ($P_ZP_Z = P_Z$)
    \item Let's define $\tilde{y} = P_Zy$, $\tilde{u} = P_Zu$ and $\tilde{X} = P_ZX$
    \item It's as if we had premultiplied the original model by $P_Z$: $\tilde{y} = \tilde{X}b + \tilde{u}$
    \item The IV estimator is: $\hat{b}_{iv} = (\tilde{X}'\tilde{X})^{-1}\tilde{X}'\tilde{y} = (X'P_ZX)^{-1}X'P_Zy$
\end{itemize}
\end{frame}

\begin{frame}{Remark}
\begin{itemize}
    \item To get to the estimated model, we multiply both $y$ and $X$ on the left by $P_Z$
    \item What would happen if we multiplied only $X$ and not $y$?
    \item $(\tilde{X}'\tilde{X})^{-1}\tilde{X}'y = (X'P_ZX)^{-1}X'P_Zy = \hat{b}_{iv}$
    \item It does not change anything as regards $\hat{b}_{iv}$: in practice, we transform only $X$ and leave $y$ unchanged
\end{itemize}
\end{frame}

\begin{frame}{Intuition (2)}
We can also explain this geometrically:
\begin{itemize}
    \item To remove the correlation between $X$ and $u$, we project the model onto the vector space $L(Z)$, spanned by the columns of $Z$, which are at the same time exogenous and correlated with $X$
    \item This amounts to keeping from $X$ only the exogenous information, uncorrelated with the error terms
    \item The more the $Z$ are correlated with the $X$ (and the more numerous the $Z$'s are), the more precise the estimator is, because the information loss is minimal
    \item $P_Zy$ is the prediction from the regression of $y$ on variables $Z$, same for $P_ZX$
\end{itemize}
\end{frame}

\begin{frame}{Remarks}
\begin{itemize}
    \item If $p = k$ (same number of instruments $Z$ and of explanatory variables $X$), we get $\hat{b}_{iv} = (Z'X)^{-1}Z'y$
    \item Proof: rewriting $\hat{b}_{iv}$, given that in this case, matrix $Z'X$ is square and invertible
    \item But with this minimal number of instruments, we cannot test the validity of the instruments (Sargan test, see below)
\end{itemize}
\end{frame}

\begin{frame}{Generalization (1)}
\begin{itemize}
    \item What if some variables in $X$ are endogenous, and others exogenous?
    \item Let's call $X_1$ the set of the exogenous ones and $X_2$ the set of the endogenous ones: the matrix notation would be $X = (X_1, X_2)$ (again, the constant is always exogenous)
    \item In that case, variables $X_1$ can be used as instruments in addition to the $Z$
    \item Let's call $W$ this extended set of instruments: $W = (Z, X_1)$
    \item To compute the IV estimator, we premultiply $X$ by $P_W$
    \item $P_WX = P_W(X_1, X_2) = (P_WX_1, P_WX_2) = (X_1, P_WX_2)$
    \item Indeed, since $X_1$ belongs to $W$, it is projected onto itself and remains unchanged
\end{itemize}
\end{frame}

\begin{frame}{Generalization (2)}
\begin{itemize}
    \item Conclusion: we will use $X_1$ as additional instruments because they will remain unchanged in the transformed model, and will help to instrument the $X_2$
    \item We only need to make sure that there are at least as many ``new'' instruments $Z$ as endogenous variables $X_2$, and don't rely only on the exogenous $X_1$
    \item $Z$ are sometimes called \emph{excluded instruments}
    \item We realize here that the constant, exogenous by nature, belongs to the exogenous $X_1$, and that we were right to use it as an instrument in the first place
    \item That's what we did when describing the first stage regression, which contains a constant by default
\end{itemize}
\end{frame}

\begin{frame}{Important: final notation}
\begin{itemize}
    \item Since the exogenous variables from the set of variables $X$ will always be used as instruments, we don't usually bother mentioning it
    \item \textbf{We usually call $Z$ the global matrix of instruments that contains both the excluded instruments and the variables from $X$ that are exogenous}
    \item So that the hypotheses and formulae are the same as the ones presented in the beginning
    \item We only need to make sure that there are at least as many excluded instruments as endogenous variables
\end{itemize}
\end{frame}

\begin{frame}{Properties of the IV estimator (1)}
\begin{itemize}
    \item $\hat{b}_{iv}$ is biased in a small sample
    \item We cannot derive a general expression for its finite-sample variance-covariance matrix
    \item But it is consistent (when the sample size goes to infinity)
    \item It is asymptotically normal:

    \[\sqrt{N}(\hat{b}_{iv} - b) \xrightarrow{d} \mathcal{N}\left(0, \sigma^2(V_{ZX}'V_Z^{-1}V_{ZX})^{-1}\right)\]
\end{itemize}
\end{frame}

\begin{frame}{Properties of the IV estimator (2)}
There is no general expression for its variance-covariance matrix, but we can derive its \emph{asymptotic} variance-covariance matrix, which is consistently estimated by:

\[\widehat{V}_{asympt}\left(\sqrt{N}(\hat{b}_{iv} - b)\right) = \hat{\sigma}_{iv}^2\left(\frac{X'P_ZX}{N}\right)^{-1}\]

with $\hat{\sigma}_{iv}^2 = \frac{SSR_{iv}}{N} = \frac{\hat{u}'\hat{u}}{N}$ and $\hat{u} = y - X\hat{b}_{iv}$

\begin{itemize}
    \item When OLS is consistent, IV is less precise than OLS
    \item Precision increases with the number of instruments (asymptotically)
\end{itemize}
\end{frame}

\begin{frame}{Tests on regression parameters}
\begin{itemize}
    \item We do not have convenient properties of the estimator in the finite sample case
    \item We can only run asymptotic tests: Wald tests, similar to F-tests, but when the sample size is large
    \item Example of a test for $r$ linear constraints on parameters:
\end{itemize}

\[f = \frac{(C\hat{b}_{iv} - Cb)'[C(X'P_ZX)^{-1}C']^{-1}(C\hat{b}_{iv} - Cb)}{\hat{\sigma}_{iv}^2} \xrightarrow{d} \chi_r^2\]

With $r$ the number of linear restrictions: $C$ is of size $(r,k)$.

Proof: taking the expression of the asymptotic normality of $\hat{b}_{iv}$, taking its quadratic form and ``dividing'' it by its variance gives a $\chi^2$ distribution (the ``true'' variance is replaced by its consistent estimator).
\end{frame}

% Section 4: IV in practice
\section{IV in practice}

\begin{frame}{How to run IV estimation}
\begin{itemize}
    \item Python: \texttt{IV2SLS} from \texttt{statsmodels.sandbox.regression.gmm} (basic), or the \texttt{linearmodels} package (\texttt{linearmodels.iv.IV2SLS}), which also reports first-stage statistics and the exogeneity and Sargan tests
    \item This amounts to running two-stage least squares
    \item Intuition: first, regress the $y$ and $X$'s on variables $Z$, then use the predictions in the model instead of the original values
\end{itemize}
\end{frame}

\begin{frame}{Two-stage least squares (1)}
\begin{itemize}
    \item Run $k+1$ regressions, to get $P_Zy$ and $P_ZX$
    \item Estimate OLS on the transformed model $P_Zy = P_ZXb + \tilde{u}$
    \item We thus get $\hat{b}_{iv} = (X'P_ZX)^{-1}X'P_Zy$
    \item The first-stage regressions can be used to assess the relevance of the instruments (they have to be correlated enough with the $X$'s)
    \item Remark: this provides the same values if we do not replace $y$ by $P_Zy$
\end{itemize}
\end{frame}

\begin{frame}{Two-stage least squares (2)}
\textbf{Warning}: if we do this procedure ``by hand'', running 2 OLS regressions, instead of running the dedicated IV procedure of the software, the standard errors of the second regression cannot be used for tests on the coefficients

\textbf{Reason}: in the second stage equation, residuals are computed as: $\hat{u} = P_Zy - P_ZX\hat{b}_{iv}$ (or $y - P_ZX\hat{b}_{iv}$)

Whereas they should be computed as $\hat{u} = y - X\hat{b}_{iv}$
\end{frame}

\begin{frame}{Exogeneity test (1)}
\begin{itemize}
    \item We test $H_0$: $E(X'u) = 0$
    \item This is called the ``Hausman test'' or ``Durbin-Wu-Hausman test''
    \item If $H_0$ is true, then both OLS and IV estimators are consistent (and OLS is more precise)
    \item If $H_0$ is false, only the IV estimator is consistent
    \item The test is based on the difference between $\hat{b}_{iv}$ and $\hat{b}_{ols}$
    \item They are asymptotically normal: if we compute the difference between the two, take its quadratic form and ``divide'' it by its variance matrix, we will get a $\chi^2$ distribution
\end{itemize}
\end{frame}

\begin{frame}{Exogeneity test (2)}
One could think that the parameter of this $\chi^2$ is the number of tested variables (those potentially endogenous), just like a Wald test. But looking closely at the test statistic:

\[H = (\hat{b}_{iv} - \hat{b}_{ols})'(V(\hat{b}_{iv}) - V(\hat{b}_{ols}))^{-1}(\hat{b}_{iv} - \hat{b}_{ols})\]

This parameter is in fact equal to the rank of the matrix $V(\hat{b}_{iv}) - V(\hat{b}_{ols})$ (whose inverse is then a generalized inverse).

In finite samples, the estimated difference of variances may not be positive definite, so sometimes the software can't run the test (or finds $H < 0$).
\end{frame}

\begin{frame}{The augmented regression}
\begin{itemize}
    \item Consider model $y = Xb + u$, where a subset $x$ of the variables belonging to $X$ might be endogenous
    \item Let's call $Z$ the instruments, some belonging to $X$ (in fact the $X$ without the $x$) and some not
    \item Consider the augmented model: $y = Xb + M_Zxc + \varepsilon$
    \item $M_Zx$ are the residuals of the regression of $x$ on $Z$
    \item The $\hat{b}$ of this ``augmented'' regression is equal to the IV estimator of the original model
    \item Testing $c = 0$ amounts to testing the exogeneity of $x$ (it is equivalent to the Hausman test)
\end{itemize}
\end{frame}

\begin{frame}{Proofs}
\begin{itemize}
    \item $\hat{b}_{aug} = \hat{b}_{iv}$ and equivalence of tests: using the Frisch-Waugh theorem
    \item Remark: this augmented model has no theoretical meaning, it's only a tool
\end{itemize}
\end{frame}

\begin{frame}{Selecting valid instruments}
Sargan test: $H_0$: $E(Z'u) = 0$

Also called: test of overidentifying restrictions

Under $H_0$:

\[\frac{\hat{u}'P_Z\hat{u}}{s^2} \xrightarrow{d} \chi_{p-k}^2\]

with $\hat{u} = y - X\hat{b}_{iv}$ and $s^2 = \frac{\hat{u}'\hat{u}}{N}$.

$\hat{u}'P_Z\hat{u}$ is the sum of the squared fitted values of the regression of $\hat{u}$ on $Z$ (the statistic is $N R^2$ of this regression).

Remark: when $p = k$, the statistic is always zero because $\hat{b}_{iv} = (Z'X)^{-1}Z'y$ and $Z'\hat{u} = 0$. So we cannot run the test with the minimum number of instruments.
\end{frame}

\begin{frame}{The problem with weak instruments}
\begin{itemize}
    \item If instruments are too weakly correlated with the $X$'s, even if we increase the number of observations, there can be an important bias in estimations
    \item Plus, the estimator has a non-normal sampling distribution, which makes statistical inference meaningless
    \item The weak instrument problem is increased with many instruments, so drop the weakest ones and use the most relevant ones
    \item To measure how instruments are correlated with the potentially endogenous variables, run the first-stage regression explaining the latter by the instruments
    \item Criterion: the F statistic testing that the \textit{excluded} instruments are jointly zero in the first stage (not the global F): if $F < 10$, instruments are weak (Staiger-Stock rule of thumb)
\end{itemize}
\end{frame}

\begin{frame}{Returns to education}
Angrist-Krueger (1991), ``Does Compulsory School Attendance Affect Schooling and Earnings?'', Quarterly Journal of Economics

\begin{itemize}
    \item Goal: find the causal link between education and wages, knowing that such a wage equation is subject to endogeneity
    \item In the US, compulsory education starts the year pupils turn 6, and ends when they turn 16
    \item Pupils born in January begin school at 6 years 8 months old; those born in December, at 5 years 9 months old
    \item For pupils who stop school at 16, those born in December study almost one year more than those born in January
    \item Quarter of birth is thus correlated with schooling, but as good as random, so uncorrelated with individual ability: good instrument for schooling
\end{itemize}
\end{frame}

\begin{frame}{Quarter of birth and education}
\begin{center}
\includegraphics[width=0.8\textwidth]{quarter_birth_education.png}
\end{center}
\end{frame}

\begin{frame}{Quarter of birth and earnings}
\begin{center}
\includegraphics[width=0.8\textwidth]{quarter_birth_earnings.png}
\end{center}
\end{frame}

\begin{frame}{Results}
\begin{center}
\includegraphics[width=0.9\textwidth]{angrist_krueger_results.png}
\end{center}
\end{frame}

\begin{frame}{Method}
\begin{itemize}
    \item Instruments: quarter-of-birth dummies, interacted with year-of-birth (and state-of-birth) dummies
    \item Result: OLS underestimates the impact of education on earnings
    \item Caveat: Bound, Jaeger and Baker (1995) showed that with so many weak instruments, 2SLS is biased towards OLS (see weak instruments)
    \item Even better method: randomized trials
    \item E.g.\ Moving to Opportunity, RAND Health Insurance Experiment
    \item Many other methods in policy evaluation: differences-in-differences, matching, etc.
\end{itemize}
\end{frame}

\begin{frame}{Are IV estimates reliable?}
\begin{itemize}
    \item The goal of IV estimation is to bring to light real causality and not mere correlation between outcome $y$ and explanatory variables $X$
    \item In treatment evaluation, we want to assess the Average Treatment Effect
    \item So we instrument the dummy indicating whether the individual received the treatment, because we suspect that taking the treatment is not exogenous: it is likely to be correlated with unobserved heterogeneity
    \item However, it can be shown that in doing so, we estimate a Local Average Treatment Effect (Imbens and Angrist, 1994), i.e.\ the impact of the treatment only on the individuals for whom the instrument indeed has an impact (called \emph{compliers})
\end{itemize}
\end{frame}

\end{document}
